% Generated by tools/edn_latex.py from reports/edn.md.
% Do not edit: change reports/edn.md and run `make edn`.
\ednentry{
  id = {EDN-40},
  title = {The \texorpdfstring{\texttt{ES}}{ES} holdout stays row-selected, so a cross-border dealer is reported rather than moved},
  date = {2026-09-30},
  milestone = {M2: Reproducibility (shapes M6: Monitoring)},
  topic = {Evaluation Protocol, Monitoring},
  participants = {@lukas2510},
  decision = {The \texttt{split} stage selects the holdout per listing, not per seller, so ``every \texttt{ES} listing is held out'' (EDN-03) takes precedence over ``no seller appears in two sets'' (EDN-14) for a dealer that lists in \texttt{ES} and elsewhere. The stage counts any such seller and reports it as a warning with row counts, the disjointness tests include the holdout instead of iterating the split names only, and the count is pinned by a test.},
  rationale = {\emph{Alternatives considered:}
\begin{itemize}[leftmargin=*, topsep=2pt]
\item \textbf{Option A: hold out whole sellers that touch \texttt{ES}.} Both invariants would then hold without exception.
\newline Pros: the holdout would be seller-disjoint from the training sets, so an \texttt{ES} replay could not contain a dealer the model trained on, which is the property EDN-14's finding says matters.
\newline Cons: it changes what the holdout \emph{is}. EDN-03 defines it as the \texttt{ES} listings and prices its cost as exactly those rows, so the artefact would gain non-\texttt{ES} rows that M6 replays as Spanish traffic, and its size would no longer be the 6,079 rows EDN-03 and the dataset card state. On the real snapshot it would also buy nothing: 0 sellers list both inside and outside \texttt{ES}, because \texttt{hash\_\allowbreak{}seller\_\allowbreak{}group} keys a dealer by its company name and no name occurs on both sides.
\item \textbf{Option B (chosen): keep the row selection, detect the conflict and report it with counts.}
\newline Pros: EDN-03's definition of the holdout is untouched, and on the real snapshot the resolution costs nothing measurable. The violation becomes visible where it does occur rather than being silently excluded from the tests: the fixtures the tests run on have 91 such sellers at 2,000 raw rows and 898 at 20,000, and no test would have noticed, because both disjointness tests iterated the split names and skipped the holdout.
\newline Cons: a weaker invariant than option A. The stage does not guarantee seller-disjointness across all five artefacts, only across the four split sets, and a future snapshot with cross-border dealer names would need this decision revisited rather than being handled automatically.
\item \textbf{Option C: leave it as it was, selected per row and unmentioned.}
\newline Cons: this is what the pull request did, and it is the reason the conflict went unnoticed. Two documented invariants cannot both hold, and nothing said which one wins or what it costs.
\end{itemize}
\par
\emph{Why:} The two rules genuinely conflict and one of them has to lose, so the only bad answer is not saying which. EDN-03 is the authority on what the holdout contains, and its whole purpose is that the replayed rows are a market: rows from other countries in it would change what the drift measurement measures, and the 6,079-row figure is already cited in the dataset card and in EDN-14's count of the \texttt{ES} rows the API would accept. Against that, option A's benefit on the real snapshot is exactly zero, because no dealer name crosses the border there.
\par
What makes option B acceptable rather than a shrug is that the cost is now measured and visible. EDN-14's finding is that dealer-level shift is as large as country-level shift, so a replay containing dealers the model trained on understates drift and confounds both NFR-11's flagging and FR-15's retrain comparison. That is a real limitation of the synthetic fixture as a drift rehearsal, and it is now a reported number with a test pinning it, so if a future snapshot does grow cross-border dealer names the stage says so on every run instead of quietly violating the invariant.},
  ai = {Information seeking, Alternative generation, Alternative assessment, Recommendation},
  response = {Accepted},
  assessment = {The adversarial review of PR \#53 found the conflict, which the pull request never mentioned, and measured that it occurs 91 and 898 times in the two fixtures the tests run on while occurring 0 times on the real snapshot -- the combination that explains why it was invisible. It also identified the mechanism, that both disjointness tests iterate the split names and therefore exclude the holdout by construction. It offered both resolutions and accepted either, provided the invariant became visible. Lukas chose the row selection, because EDN-03 owns the definition of the holdout and option A would have changed a figure already published for no gain on the real data.},
  aievidence = {Claude Code adversarial review of PR \#53 on 2026-09-30, which measured the cross-border counts in both fixtures and on the real snapshot and named the conflict between EDN-03 and EDN-14.},
  evidence = {\href{https://github.com/mlops-2627q1-mds-upc/MLOPS_Recommenditos/blob/main/recommenditos/data/split_data.py}{\texttt{recommenditos/\allowbreak{}data/\allowbreak{}split\_\allowbreak{}data.\allowbreak{}py}} (\texttt{cross\_\allowbreak{}holdout\_\allowbreak{}sellers}); \href{https://github.com/mlops-2627q1-mds-upc/MLOPS_Recommenditos/blob/main/tests/test_split.py}{\texttt{tests/\allowbreak{}test\_\allowbreak{}split.\allowbreak{}py}}; \href{https://github.com/mlops-2627q1-mds-upc/MLOPS_Recommenditos/blob/main/reports/analysis/split_gate.py}{\texttt{reports/\allowbreak{}analysis/\allowbreak{}split\_\allowbreak{}gate.\allowbreak{}py}}; \hyperref[EDN-03]{EDN-03}; \hyperref[EDN-14]{EDN-14}; \href{https://github.com/mlops-2627q1-mds-upc/MLOPS_Recommenditos/issues/35}{issue \#35}; \href{https://github.com/mlops-2627q1-mds-upc/MLOPS_Recommenditos/pull/53}{PR \#53}.},
}
